\documentclass[a4paper, oneside, 12pt, titlepage]{article}

\usepackage[utf8]{inputenc}
\usepackage[T1]{fontenc}
\usepackage[main=canadien, english]{babel}
\usepackage{float}
\usepackage[inline]{enumitem}
\usepackage{pdflscape}
\usepackage{geometry}
\usepackage{array}
\usepackage[pdftex]{graphicx}

\newcolumntype{L}[1]{>{\raggedright\let\newline\\\arraybackslash\hspace{0pt}}m{#1}}
\newcolumntype{C}[1]{>{\centering\let\newline\\\arraybackslash\hspace{0pt}}m{#1}}
\newcolumntype{R}[1]{>{\raggedleft\let\newline\\\arraybackslash\hspace{0pt}}m{#1}}
\newcommand{\HRule}{\rule{\linewidth}{0.5mm}}

\floatstyle{boxed}
\restylefloat{figure}

\begin{document}

% \title{Formalisation de systèmes de types à l'aide d'Isabelle/HOL}
% \author{Martin Desharnais}
% \maketitle

\begin{titlepage}
  \begin{center}
    % Upper part of the page. The '~' is needed because \\
    % only works if a paragraph has started.
    \includegraphics[width=6cm]{./logo-ETS.png}~\\[1cm]

    \textsc{\LARGE École de technologie supérieure}\\[1.5cm]

    \textsc{\Large LOG792 Projet de fin d'étude}\\[0.25cm]
    \textsc{\Large Rapport d'étape}\\[0.5cm]

    % Title
    \HRule \\[0.4cm]
    { \huge \bfseries Formalisation de systèmes de types à l'aide d'Isabelle/HOL \\[0.4cm] }

    \HRule \\[1.5cm]

    % Author and supervisor
    \begin{minipage}{0.4\textwidth}
    \begin{flushleft} \large
    \emph{Auteur :}\\
    Martin Desharnais
    \end{flushleft}
    \end{minipage}
    \begin{minipage}{0.4\textwidth}
    \begin{flushright} \large
    \emph{Superviseur :} \\
    Dr.~David Labbé
    \end{flushright}
    \end{minipage}

    \vfill

    % Bottom of the page
    {\large \today}

  \end{center}
\end{titlepage}

\pagestyle{empty}
\pagenumbering{Roman}
\tableofcontents
%\clearpage
%\listoftables
%\clearpage
%\listoffigures

\cleardoublepage
\pagenumbering{arabic}

\section{Problématique et contexte}

% \emph{Quel sont les problèmes généraux et spécifiques que vous cherchez à solutionner – totalement
% ou en partie? Quel est le contexte de ces problèmes? Quelle est la situation actuelle?}

Ce projet s'intéresse à l'étude des systèmes de types, dans le contexte des langages de
programmation, dont voici une définition \cite{pierce2002} :

\begin{quotation}
  Un système de types est une méthode syntaxique tractable pour prouver l'absence de certains
  comportements des programmes par la classification des phrases selon le genre de valeurs qu'elles
  calculent.
\end{quotation}

L'objectif est donc de garantir, sans l'exécuter, qu'un programme est exempt de certaines erreurs
telles qu'une faute typographique dans un nom de variable, l'appel d'une fonction non supportée dans
un certain contexte ou encore une tentative de diviser un nombre par une chaine de caractères. De
tels exemples peuvent sembler simplistes, mais sont très fréquents et peuvent avoir des conséquences
désastreuses : une incohérence informatique entre les systèmes d'unités métrique et anglo-saxon a
provoqué la destruction du Mars Climate Orbiter en 1999. Bien sûr, tous les défauts ne peuvent pas
être décelés par un système de types. Cependant, il en existe un très grand nombre, de divers
niveaux d'expressivité et de complexité, qui permettent de détecter un éventail varié d'erreurs.

Lors du développement d'un nouveau système de type, un ensemble de preuves est réalisé afin de
démontrer que celui-ci respecte ses objectifs. L'étude de ces systèmes ainsi que des
preuves qui les accompagnent est le sujet du présent projet.

\section{Objectifs du projet}

% \emph{Dans cette partie, vous décrivez précisément  ce que vous allez accomplir dans le projet.
% Quels sont les résultats attendus de votre projet? Quelles sont les retombées du projet?
% Assurez-vous que vos objectifs sont "testables", c'est-à-dire que vous serez capable de démontrer
% de façon convaincante, à la fin du projet, que les objectifs ont été atteints. Assurez-vous aussi
% que vos objectifs sont en lien avec la problématique énoncée plus haut.}

% Learning (Paper-Proof of Type Systems properties & Isabelle/HOL)
% Validate Paper-Proofs
% Clarification (Paper-Proof & Isabelle/HOL)

Les objectifs de ce projet sont quadruples :
\begin{enumerate*}[label=\arabic*)]
  \item s'initier à la formalisation avec l'assistant de preuve Isabelle/HOL;
  \item s'initier à la théorie des types;
  \item valider les preuves manuelles existantes et
  \item clarifier les cas limites de ces dernières.
\end{enumerate*}

\subsection{S'initier à la formalisation Isabelle/HOL}

Il existe deux grandes catégories d'outils pour effectuer une formalisation : les prouveurs
automatiques et les assistants de preuve interactifs. Quel que soit l'outil utilisé, il faut définir
un contexte de travail et une équation que l'on veut prouver. La différence est que, dans le premier
cas, l'outil tente de trouver une preuve entièrement automatiquement alors que, dans le second
cas, il faut travailler interactivement avec l'outil pour prouver le théorème.

Isabelle/HOL est un assistant de preuve interactif utilisant la logique d'ordre supérieure. Il
permet de spécifier des formules mathématiques, algorithmes et objets dans un langage déclaratif,
fonctionnel et typé. Il est alors possible de spécifier des propriétés sur l'interaction entre les
divers éléments. Une fois le système formalisé, il est possible d'en extraire du code exécutable
correspondant aux spécifications. Ce projet est l'occasion de s'initier à la définition d'un
système formel et aux preuves interactives à l'aide d'Isabelle/HOL.

\subsection{S'initier à la théorie des types}

Plusieurs des langages de programmation dominant actuellement ont une étape de validation des types
appliquée à la compilation. La majorité des programmeurs sont donc familiarisés avec le concept.
Malheureusement, les systèmes de types présents dans ces langages sont généralement simples,
imposent nombre de contraintes à leurs utilisateurs et n'offrent qu'un nombre limité de garanties en
retour. Certaines de ces limitations ont été mitigées dans de nouvelles versions du langage (e.g.
l'ajout des types et fonctions génériques en Java).

Cependant, des options plus expressives et plus puissantes sont bien connues, ou bien
actuellement en développement par les acteurs du milieu. Ce projet est l'occasion de consolider
les acquis fondamentaux et d'en apprendre plus sur les concepts plus avancés de la théorie des
types.

\subsection{Valider les preuves manuelles existantes}

La théorie des types étant un sujet de recherche très actif depuis plusieurs dizaines d'années, un
grand nombre de publications décrivent les caractéristiques de différents systèmes de type.
Cependant, une preuve manuelle étant validée par des êtres humains, il est toujours possible que des
erreurs s'y soient glissées. La formalisation de celles-ci à l'aide d'un assistant de preuve permet
de valider, sous réserve que l'assistant de preuve soit correct, qu'aucune erreur logique n'est
présente. S'il s'avérait qu'une erreur soit découverte dans le cadre de ce projet, l'information
serait transmise à l'auteur initial afin de l'informer de la situation.

\subsection{Clarifier les cas limites des preuves manuelles}

Les propriétés énoncées et prouvées manuellement semblent souvent évidentes dès lors qu'elles sont
appliquées à un exemple concret. Cette méthode de visualisation a toutefois ses limites puisque
certaines constructions plus complexes peuvent entrainer des résultats inattendus. La formalisation
de ces propriétés à l'aide d'un assistant de preuve oblige son auteur à considérer la liste
exhaustive des constructions du langage et permet ainsi d'acquérir une meilleure compréhension de la
propriété et des cas limites. Cette technique fut utilisée avec succès afin d'enseigner l'ingénierie
logicielle à des étudiants de premier cycle \cite{page:thm-proving-undergraduates, pierce:LambdaTA,
nipkow:teaching-proof-assistant}.

\section{Méthodologie}

% \emph{Expliquez comment vous allez atteindre les objectifs du projet. Ce sont les activités
% (analyse, conception, mesure, tests, gestion de la configuration, revues, etc.) et les
% responsabilités identifiées dans cette section qui guideront l’affectation des responsabilités aux
% membres de l’équipe et aussi à l’établissement de votre échéancier de travail.}
%
% \emph{Note : une approche itérative est recommandée.}

L'ouvrage de référence de ce projet est le livre \emph{Types and Programming Languages} de
Benjamin~C.~Pierce. Ce dernier est composé de six sections : les systèmes non typés, les types
simples, le sous-typage, les types récursifs, le polymorphisme et les systèmes d'ordre supérieur.
Chaque section est composée de plusieurs chapitres décrivant un système de type bonifiant un système
décrit précédemment en lui adjoignant un concept supplémentaire. Les figures
\ref{fig:TAPL-section-1} et \ref{fig:TAPL-section-2} présentent les chapitres des deux premières
sections sur lesquelles se concentre ce projet --- ceux en gras sont ceux qui seront formalisés
--- et la figure \ref{fig:TAPL-chapter-dependencies} présente les dépendances entre ceux-ci.

Ces chapitres furent choisis, car ils décrivent des langages et leurs théorèmes au lieu d'en
expliquer la théorie ou d'en présenter une implémentation. De plus, ils culminent au
$\lambda$-calcul simplement typé, qui a la propriété de pouvoir représenter la majorité des langages
de programmation existants\footnote{Ceci n'est pas tout à fait exact puisqu'il faudrait y ajouter la
fonctionnalité de pouvoir communiquer avec l'environnement d'exécution afin de pouvoir faire des
opérations telles que lire ou écrire dans un fichier, effectuer de la communication interprocessus,
etc.}.

\begin{figure}[h]
  \begin{enumerate}[label=§ \arabic*]
      \setcounter{enumi}{2}
    \item \textbf{Expressions arithmétiques non typées}
    \item Une implémentation en ML des expressions arithmétiques
    \item \textbf{Le lambda-calcul non typé}
    \item Représentation non nommée des termes
    \item Une implémentation en ML du lambda-calcul
  \end{enumerate}
  \caption{Section I du livre de référence --- Les systèmes non typés}
  \label{fig:TAPL-section-1}
\end{figure}

\begin{figure}[h]
  \begin{enumerate}[label=§ \arabic*]
      \setcounter{enumi}{7}
    \item \textbf{Expressions arithmétiques typées}
    \item \textbf{Le lambda-calcul simplement typé}
    \item Une implémentation en ML des types simples
    \item Extensions simples
    \item Normalisation
    \item Références
    \item Exceptions
  \end{enumerate}
  \caption{Section II du livre de référence --- Les types simples}
  \label{fig:TAPL-section-2}
\end{figure}

\begin{figure}[h]
  \begin{center}
    \setlength{\unitlength}{1cm}
    \begin{picture}(4,4.5)
      \thicklines
      \put(0,0){\framebox(1,1){3}}
      \put(0,3){\framebox(1,1){8}}
      \put(0.5,3){\vector(0,-1){2}}
      \put(3,0){\framebox(1,1){5}}
      \put(3,0.5){\vector(-1,0){2}}
      \put(3,3){\framebox(1,1){9}}
      \put(3,3.5){\vector(-1,0){2}}
      \put(3.5,3){\vector(0,-1){2}}
    \end{picture}
  \end{center}
  \caption{Dépendances entre les chapitres}
  \label{fig:TAPL-chapter-dependencies}
\end{figure}

Le projet consiste à formaliser séquentiellement les différents chapitres en se basant sur le
travail des chapitres précédents. Pour cette raison, le premier chapitre fut le plus long à
formaliser : tout devant être fait à partir de zéro. Pour les chapitres suivants, la première étape
est de copier la formalisation des dépendances et de la modifier pour inclure les nouveaux concepts
introduits.

Chacune des formalisations se fait en quatre étapes :

\begin{enumerate}
  \item Lecture attentive du chapitre et compréhension générale des concepts énoncés;
  \item Définition dans Isabelle/HOL des structures nécessaires à la formalisation;
  \item Preuve des différents lemmes et théorèmes ainsi que, optionnellement, des exercices;
  \item Simplification des définitions et preuves.
\end{enumerate}

L'objectif principal du projet étant de formaliser le chapitre 9 et ce genre de formalisation étant
un projet ambitieux, un certain nombre d'actions pourraient être entreprises si le respect des
échéanciers s'avère menacé :

\begin{enumerate}
  \label{lower-objectives}
  \item Ne pas prendre en compte les exercices proposés;
  \item Définir certains lemmes comme des axiomes au lieu de les prouver;
  \item Sauter les chapitres 5 ou 8 afin de passer directement au chapitre 9.
\end{enumerate}

\subsection{Techniques et outils}

\begin{description}
  \item[Isabelle] Système générique pour l'implémentation de formalismes logiques.
  \item[Isabelle/HOL] Spécialisation d'Isabelle pour la logique d'ordre supérieur
    (\emph{Higher-Order Logic} en anglais).
  \item[Isabelle/Isar] Langage structuré permettant d'écrire des preuves plus lisibles.
  \item[Isabelle/jEdit] Environement de développement intégré pour Isabelle basé sur l'éditeur
    jEdit.
  \item[Sledgehammer] Outil appliquant des trouveurs automatiques de théorèmes ainsi que des
    solveurs de « satisfaisabilité modulo théories ».
\end{description}

\section{Livrables et planification}

\subsection{Description des artéfacts}

% \emph{Cette section identifie les artefacts qui seront produits durant le projet, ainsi que la
% planification de leur réalisation.}

\begin{description}
  \item[Fiche de renseignement]
    Formulaire fournissant le titre du projet, un cours résumé ainsi que les noms des étudiants
    impliqués et du professeur superviseur.
  \item[Proposition de projet]
    Document présentant la problématique du projet, les objectifs, la méthodologie, les livrables,
    le plan de travail, les risques ainsi que les techniques et outils utilisés.
  \item[Rapport d'étape]
    Document présentant une version étoffée et mise à jour de la proposition de projet, ainsi qu'une
    version partielle du rapport technique.
  \item[Diapositives de la présentation orale]
    Diapositives utilisées pour la présentation orale finale du projet.
  \item[Rapport final]
    Document présentant la problématique du projet, les objectifs, la méthodologie employée, les
    hypothèses, les résultats, l'analyse des résultats, les conclusions, les recommandations et les
    références.
  \item[Théories Isabelle/HOL]
    Fichiers sources utilisés par l'assistant de preuve Isabelle/HOL pour sauvegarder les
    définitions et théorèmes formalisés au cours de ce projet.
\end{description}

\subsection{Planification}

% Voir Annexe A. Commentez le tableau de l’annexe A. Une approche itérative est recommandée.

Le plan de projet se trouve à l'annexe A.

\section{Risques}

% \emph{Vous devez énumérer les difficultés probables que vous pourriez rencontrer et qui pourraient
% avoir un impact sur la réalisation de votre projet. Vous devez expliquer dans un tableau chaque
% risque identifié, son impact ainsi que les moyens que vous allez mettre en œuvre pour le gérer et
% atténuer sa probabilité ou son impact.}
%
% \emph{La forme d'un risque devrait être négative. C'est-à-dire, un risque est un événement que l'on
% veut éviter. Par exemple "expérience du client" n'est pas un risque tandis que "manque
% d’expérience du client" en serait un. Voici un exemple de risque : "Un passager est blessé dans
% une voiture lors d'une collision." Une mitigation à un risque est la stratégie qui vise à éviter
% que l'événement négatif se produise. Par exemple : "Le passager porte une ceinture de sécurité et
% la voiture est équipée de coussin de sécurité gonflable." Soyez spécifique à votre projet. "Manque
% de temps pour finir le projet" et "Portée trop ambitieuse" sont des risques pour n'importe quel
% projet. Il n’est pas mauvais d’avoir des risques génériques, mais il est important d’aussi trouver
% plusieurs risques spécifiques à votre projet, de même que de trouver des mitigations qui sont
% spécifiques à votre projet.}
%
% \emph{Si vous manquez d'inspiration, vous pouvez consulter la liste LOG\_GTI\_792\_Generic Software
% Risk Factors.xls, disponible sur le site Web du cours sous la rubrique ‘Gabarits et guides’.
% Toutefois, faire attention de ne pas ainsi identifier des risques et mitigations qui sont
% applicables à la majorité des projets.}

Un certain nombre de risques sont identifiés dans le tableau \ref{tab:risks}. En comparaison de la
proposition de projet, les modifications suivantes ont été apportées :

\begin{itemize}
  \item Ajout d'une évaluation subjective des impacts.
  \item Ajout d'une explications pour les probabilités.
  \item Diminution de la probabilité du risque 2 de « Forte » à « Moyenne ».
  \item Diminution de la probabilité du risque 3 de « Forte » à « Moyenne ».
\end{itemize}

\begin{table}[!h]
  \makebox[\textwidth][c]{
  \small
  \begin{tabular}{|p{3.5cm}|p{3.5cm}|p{3.59cm}|p{3.5cm}|}
    \hline
    \textbf{Risque} &
    \textbf{Impact} &
    \textbf{Probabilité} &
    \textbf{Mitigation / \newline atténuation} \\
    \hline
    Manque d'expérience avec la théorie des types &
    Faible : Monopolise du temps pour apprendre la théorie. &
    Faible : Le $\lambda$-calcul simplement typé est le plus simple de sa catégorie. &
    Étudier attentivement l'ouvrage de référence. \\
    \hline
    Manque d'expérience avec l'assistant de preuve Isabelle/HOL &
    Moyen : Monopolise du temps pour apprendre le fonctionnement de l'outil. &
    Moyenne : La majorité de l'apprentissage a été effectuée lors de la formalisation du premier
    chapitre. &
    Étudier attentivement et se référer au besoin à la documentation de l'outil. \\
    \hline
    Manque d'expérience en formalisation &
    Fort : Monopolise du temps pour apprendre la méthodologie. &
    Moyenne : La majorité de l'apprentissage a été effectuée lors de la formalisation du premier
    chapitre. &
    S'informer auprès des chercheurs expérimentés de la chaire de recherche. \\
    \hline
    Objectifs trop ambitieux &
    Fort : Ne pas réaliser toutes les formalisations prévues. &
    Moyenne : La formalisation des chapitres s'est fait selon l'échéancier prévu. &
    Un certain nombre de mitigations sont décrites à la page \pageref{lower-objectives} de la
    section \ref{lower-objectives}. \\
    \hline
  \end{tabular}
  }
  \caption{Risques et mitigations}
  \label{tab:risks}
\end{table}

\section{Références}

% \emph{Vos références initiales. Par exemple: documents de projet existants, compilation de patrons,
% documentation sur la technologie, etc.}

\renewcommand*{\refname}{\vspace*{-1em}}\vspace*{-1em}

\nocite{*}
\bibliographystyle{alpha}
\bibliography{bibliography}

\section{Table des matières du rapport}

\emph{Sur recommendation du directeur du département de génie logiciel et des TI, monsieur Jean-Marc
Robert, le doyen des études, monsieur Pierre Bourque, a autorisé que la rédaction du rapport soit
effectuée en anglais.}

\begin{otherlanguage}{english}
  This section discusses a possible thesis outline. It consists of seven main sections: three to
  introduce reader to the formalization of type systems and four containing the proper formlizations.

  \begin{enumerate}[label=\textbf{\Roman*}]
    \item \textbf{Introduction}
      \begin{enumerate}[label=\textbf{\arabic*}]
        \item \textbf{Introduction to Type Systems} ~ \\
          This section will provide a brief introduction to type systems, their history, current use
          and possible future directions.
        \item \textbf{Introduction to Formalization} ~ \\
          This section will provide a brief introduction to logic, formal reasoning and formalization.
        \item \textbf{Introduction to Isabelle} ~ \\
          This section will provide a brief introduction to the Isabelle/HOL: a proof
          assistant based on higher-order logic.
      \end{enumerate}
    \item \textbf{Untyped Systems}
      \begin{enumerate}[label=\textbf{\arabic*}]
        \setcounter{enumii}{3}
        \item \textbf{Untyped Arithmetic Expression} ~ \\
          This section will present the formalization of the language introduced in the Chapter 3 of
          the reference book. It will consist of the definitions and theorems defining the formal
          system.
        \item \textbf{Untyped Lambda Calculus} ~ \\
          This section will present the formalization of the language introduced in the Chapter 5 of
          the reference book. It will consist of the definitions and theorems defining the formal
          system. It will also feature a brief introduction to the $\lambda$-calculus and its relation
          to programming language theory.
      \end{enumerate}
    \item \textbf{Simple Types}
      \begin{enumerate}[label=\textbf{\arabic*}]
        \setcounter{enumii}{5}
        \item \textbf{Typed Arithmetic Expressions} ~ \\
          This section will present the formalization of the language introduced in the Chapter 8 of
          the reference book. It will consist of the definitions and theorems defining the formal
          system.
        \item \textbf{Simply Typed Lambda Calulus} ~ \\
          This section will present the formalization of the language introduced in the Chapter 9 of
          the reference book. It will consist of the definitions and theorems defining the formal
          system.
      \end{enumerate}
    \item \textbf{conclusion}
  \end{enumerate}
\end{otherlanguage}

\newgeometry{margin=2.5cm}
\begin{landscape}
\thispagestyle{empty}

\section{Annexe A : Plan de travail}

% \emph{Le tableau suivant présente la planification pour la réalisation des tâches ou artefacts
% décrits précédemment.}
%
% \emph{Note : il faut compléter la liste des tâches et des jalons. Assurez-vous que vos tâches sont
% exprimées en termes d'activités d'ingénierie, par exemple analyse, architecture, conception,
% implémentation, etc, et non en termes de production d'artéfacts (ex: rédaction document X). Les
% artéfacts affectés par une activité d'ingénierie sont indiqués dans la colonne
% "Livrables/Artéfacts". Si vous utilisez une approche itérative, rendez vos itérations visibles
% dans ce plan. Si le projet est réalisé en équipe, ajoutez une colonne d'effort pour chaque membre
% de l'équipe.}

\begin{table}[!h]
  \small
  \centering
  \begin{tabular}{|l||l|l|p{1.6cm}|p{1.6cm}||l|l|}
    \hline
    \textbf{\#} &
    \textbf{Commence} &
    \textbf{Termine} &
    \textbf{Efforts estimés (heures)} &
    \textbf{Efforts actuels (heures)} &
    \textbf{Jalon} &
    \textbf{Artéfacts} \\
    \hline \hline
    1   & & 2014-09-02 & 1 & 1 & Remise de la fiche de renseignement & Fiche de renseignements \\
    1.1 & & 2014-09-01 & 1 & 1 & Rencontre --- professeur superviseur & \\
    \hline
    2   & 2014-09-02 & 2014-09-19 & 5 & 5 & Remise de la proposition de projet & Proposition de projet \\
    2.1 &            & courriels  & 1 & 1 & Rencontre --- professeur superviseur & \\
    \hline
    2.2.1 & 2014-09-19 &            & 2  & 2  & Étude du chapitre 3 & \\
    2.2.2 &            &            & 10 & 8  & Définitions du chapitre 3 & \\
    2.2.3 &            &            & 15 & 25 & Preuves du chapitre 3 & \\
    2.2.4 &            & 2014-10-10 & 2  & 2  & Simplification du chapitre 3 & Théorie Isabelle/HOL \\
    \hline
    3   & 2014-09-19 & 2014-10-24   & 10 & & Remise du rapport d'étape & Rapport d'étape \\
    3.1 &            & à déterminer & 1  & & Rencontre --- professeur superviseur & \\
    \hline
    3.2.1 & 2014-10-10 &            & 2  & 2  & Étude du chapitre 5 & \\
    3.2.2 &            &            & 5  & 2  & Définitions du chapitre 5 & \\
    3.2.3 &            &            & 10 & 10 & Preuves du chapitre 5 & \\
    3.2.4 &            & 2014-10-24 & 3  & 2  & Simplification du chapitre 5 & Théorie Isabelle/HOL \\
    \hline
    3.3.1 & \textbf{2014-11-03} &                     & 2  & & Étude du chapitre 8 & \\
    3.3.2 &                     &                     & 5  & & Définitions du chapitre 8 & \\
    3.3.3 &                     &                     & 10 & & Preuves du chapitre 8 & \\
    3.3.4 &                     & \textbf{2014-11-10} & 3  & & Simplification du chapitre 8 & Théorie Isabelle/HOL \\
    \hline
    3.4.1 & \textbf{2014-11-10} &                     & 2  & & Étude du chapitre 9 & \\
    3.4.2 &                     &                     & 5  & & Définitions du chapitre 9 & \\
    3.4.3 &                     &                     & 15 & & Preuves du chapitre 9 & \\
    3.4.4 &                     & \textbf{2014-11-21} & 3  & & Simplification du chapitre 9 & Théorie Isabelle/HOL \\
    \hline
    5 & \textbf{2014-11-25} & 2014-12 & 10 & & Présentation & Diapositives de la présentation orale \\
    \hline
    6   & 2014-12-01 & 2014-12      & 30 & & Remise du travail & Rapport final \\
    6.1 &            & à déterminer & 1  & & Rencontre --- professeur superviseur & \\
    \hline
  \end{tabular}
\end{table}

\end{landscape}
\restoregeometry

\end{document}
